\documentclass[9pt]{extarticle}
\usepackage[utf8]{inputenc}
\usepackage[T1]{fontenc}
\usepackage[a4paper,margin=1.1cm,top=1.4cm,bottom=1.1cm]{geometry}
\usepackage{amsmath,amssymb,mathtools}
\usepackage{multicol}
\usepackage{xcolor}
\usepackage{fancyhdr}
\definecolor{acento}{HTML}{0A6ED6}
\setlength{\columnsep}{0.9cm}
\setlength{\columnseprule}{0.3pt}
\setlength{\parindent}{0pt}
\setlength{\abovedisplayskip}{4pt}\setlength{\belowdisplayskip}{4pt}\setlength{\abovedisplayshortskip}{2pt}\setlength{\belowdisplayshortskip}{4pt}
\pagestyle{fancy}\fancyhf{}
\renewcommand{\headrulewidth}{0.4pt}
\fancyhead[L]{\small\textbf{Formulario} \textperiodcentered\ EIE2211 Electromagnetismo PUCV \textperiodcentered\ Teoría de circuitos}
\fancyhead[R]{\small Pág. \thepage}
\newcommand{\grupo}[1]{\vspace{5pt}{\color{acento}\small\bfseries\MakeUppercase{#1}}\par\vspace{1pt}{\color{acento}\hrule height 0.5pt}\vspace{3pt}}
\newcommand{\formula}[3]{\par\vspace{2pt}\begin{minipage}{\linewidth}#1\textbf{#2}\par\vspace{-3pt}
  \[ \begin{gathered} #3 \end{gathered} \]\vspace{-6pt}\end{minipage}\par}
\begin{document}
\begin{multicols}{2}
\formula{\grupo{Carga y corriente}}{Principio de continuidad}{I = \oint_S \vec J\cdot d\vec s = -\frac{dQ_i}{dt}\\[2pt]\nabla\cdot\vec J = -\frac{\partial\rho}{\partial t}}
\formula{}{Corriente y densidad de corriente}{I = \int_S \vec J\cdot d\vec s  \;\Longrightarrow\;  J = \frac{I}{S}\ \ (\vec J\ \text{uniforme y perpendicular a } S)}
\formula{\grupo{Conductores y resistencia}}{Velocidad de deriva y movilidad}{\vec U_d = -\mu_e\,\vec E\quad [\mathrm{m/s}]}
\formula{}{Densidad de corriente y ley de Ohm puntual}{\vec J = \rho_e\,\vec U_d\\[2pt]\sigma = -\rho_e\,\mu_e\\[2pt]\vec J = \sigma\,\vec E}
\formula{}{Resistencia de un conductor uniforme}{R = \frac{L}{\sigma S} = \rho_\Omega\,\frac{L}{S}\quad [\Omega]}
\formula{}{Ley de Ohm en circuitos}{V_{ab} = R\,I\\[2pt]I = G\,V_{ab}, \quad G = \frac{1}{R}\ [\mathrm{S}]}
\formula{}{Conductividad de un semiconductor}{\sigma_{sc} = -\rho_e\mu_e + \rho_h\mu_h}
\formula{\grupo{Potencia y energía}}{Potencia con referencia pasiva}{p(t) = v(t)\,i(t)\\[2pt]p > 0:\ \text{recibe energía}\\[2pt]p < 0:\ \text{entrega energía}}
\formula{}{Potencia disipada en un resistor (efecto Joule)}{P = V I = I^2 R = \frac{V^2}{R}\\[2pt]W[t_0,t] = \int_{t_0}^{t} p(t')\,dt'}
\formula{\grupo{Fuentes}}{Fuerza electromotriz}{\varepsilon = \frac{W_{\text{no est}}}{q} = \int_{-}^{+} \vec f_s\cdot d\vec l\quad [\mathrm{V}]}
\formula{}{Tensión en bornes de una fuente real}{V_{ab} = \varepsilon - I\,r\ \ \text{(entrega)}\\[2pt]V_{ab} = \varepsilon + I\,r\ \ \text{(se carga)}\\[2pt]I = \frac{\varepsilon}{r + R}}
\formula{}{Potencia en la carga y máxima transferencia}{P_R = \frac{\varepsilon^2 R}{(r + R)^2}\\[2pt]P_{R,\max} = \frac{\varepsilon^2}{4r}\ \ \text{en}\ R = r}
\formula{\grupo{Leyes de Kirchhoff}}{Ley de corrientes de Kirchhoff (LCK)}{\sum I_{\text{entran}} = \sum I_{\text{salen}}  \;\Longleftrightarrow\;  \sum_k I_k = 0}
\formula{}{Divisor de corriente}{I_1 = I\,\frac{R_2}{R_1 + R_2}\\[2pt]I_2 = I\,\frac{R_1}{R_1 + R_2}}
\formula{}{Ley de voltajes de Kirchhoff (LVK)}{\oint \vec E\cdot d\vec l = 0 \;\Longrightarrow\; \sum_{\text{lazo}} \Delta V = 0\\[2pt]V_c - V_a = \sum_{a\to c}\Delta V}
\formula{}{Divisor de tensión}{V_k = V\,\frac{R_k}{\sum R}}
\formula{\grupo{Redes equivalentes}}{Resistencias en serie y en paralelo}{R_{eq} = R_1 + R_2 + \dots\\[2pt]\frac{1}{R_{eq}} = \frac{1}{R_1} + \frac{1}{R_2} + \dots\\[2pt]R_1 \parallel R_2 = \frac{R_1R_2}{R_1 + R_2}}
\formula{}{Transformación triángulo a estrella ($\Delta \to Y$)}{R_a = \frac{R_{ab}R_{ca}}{\Sigma} \quad R_b = \frac{R_{ab}R_{bc}}{\Sigma} \quad R_c = \frac{R_{bc}R_{ca}}{\Sigma}\\[2pt]\Sigma = R_{ab} + R_{bc} + R_{ca}}
\formula{}{Transformación estrella a triángulo (Y $\to$ $\Delta$)}{\Pi = R_aR_b + R_bR_c + R_cR_a\\[2pt]R_{ab} = \frac{\Pi}{R_c} \quad R_{bc} = \frac{\Pi}{R_a} \quad R_{ca} = \frac{\Pi}{R_b}}
\formula{\grupo{Métodos de análisis}}{Análisis de mallas por inspección}{\begin{bmatrix} R_{11} & -R_{12} \\ -R_{21} & R_{22} \end{bmatrix}\begin{bmatrix} I_1 \\ I_2 \end{bmatrix} = \begin{bmatrix} \sum \varepsilon_1 \\ \sum \varepsilon_2 \end{bmatrix}}
\formula{}{Supermalla}{I_2 - I_1 = I_s\\[2pt]\text{(fuente de corriente en la rama compartida)}\\[2pt]+\ \text{LVK por el contorno de ambas mallas}}
\formula{}{Análisis de nodos}{\sum_j \frac{V_i - V_j}{R_{ij}} = I_{s,\text{entrando}}  \;\Longleftrightarrow\;  \sum_j G_{ij}(V_i - V_j) = I_{s,\text{entrando}}}
\formula{}{Supernodo}{V_1 - V_2 = \varepsilon\\[2pt]+\ \text{LCK de la superficie que encierra 1 y 2}}
\formula{\grupo{Teoremas de redes}}{Superposición}{x = x' + x'' + \dots\\[2pt]x^{(k)}:\ \text{respuesta con solo la fuente } k \text{ encendida}}
\formula{}{Equivalentes de Thévenin y Norton}{V_{th} = V_{ab}\big|_{\text{abierto}}\\[2pt]I_N = I_{ab}\big|_{\text{corto}}\\[2pt]R_{th} = \frac{V_{th}}{I_N} = R_{ab}\big|_{\text{fuentes indep. apagadas}}}
\formula{}{Máxima transferencia de potencia}{R_L = R_{th}\\[2pt]P_{\max} = \frac{V_{th}^2}{4R_{th}}}
\formula{\grupo{Transitorios}}{Relación tensión-corriente del capacitor}{i_C = C\,\frac{dv_C}{dt}\\[2pt]v_C(0^+) = v_C(0^-)}
\formula{}{Respuesta de primer orden (RC y RL)}{x(t) = x(\infty) + \left[x(0^+) - x(\infty)\right]e^{-t/\tau}\\[2pt]\tau = R_{th}C \quad\text{o}\quad \tau = \frac{L}{R_{th}}}
\formula{}{Energía almacenada en un capacitor}{W = \tfrac12 C v_C^2 = \tfrac12 Q v_C = \frac{Q^2}{2C}}
\formula{}{Relación tensión-corriente del inductor}{v_L = \frac{d\lambda}{dt} = L\,\frac{di_L}{dt}, \quad \lambda = N\Phi = L\,i_L\\[2pt]i_L(0^+) = i_L(0^-)}
\formula{}{Transitorio RL}{i_L(t) = i_L(\infty) + \left[i_L(0^+) - i_L(\infty)\right]e^{-t/\tau}\\[2pt]\tau = \frac{L}{R_{th}}}
\formula{}{Energía almacenada en un inductor}{W = \tfrac12 L\,i_L^2\\[2pt]w_m = \frac{B^2}{2\mu}}
\formula{}{Carga y descarga RC}{\text{Carga: } v_C(t) = V\left(1 - e^{-t/RC}\right),\ \ i(t) = \frac{V}{R}e^{-t/RC}\\[2pt]\text{Descarga: } v_C(t) = V_0\,e^{-t/RC}}
\formula{\grupo{Resistencia en geometrías no uniformes}}{Procedimiento general para calcular una resistencia}{I \;\to\; J = \frac{I}{S(r)} \;\to\; E = \frac{J}{\sigma} \;\to\; V = \int E\,dl \;\to\; R = \frac{V}{I}\\[2pt]R = \frac{\int \vec E\cdot d\vec l}{\int_S \sigma\vec E\cdot d\vec s}}
\formula{}{Resistencia de fuga de un coaxial}{R = \frac{\ln(b/a)}{2\pi\sigma L}\\[2pt]E(r) = \frac{V}{r\,\ln(b/a)}}
\formula{}{Resistencia entre esferas concéntricas}{R = \frac{1}{4\pi\sigma}\left(\frac{1}{a} - \frac{1}{b}\right)}
\formula{}{Electrodo semiesférico (puesta a tierra)}{R = \frac{1}{2\pi\sigma}\left(\frac{1}{a} - \frac{1}{b}\right) \xrightarrow{\,b\to\infty\,} \frac{1}{2\pi\sigma a}\\[2pt]V(r) = \frac{I}{2\pi\sigma r}}
\formula{\grupo{Capacitores y dieléctricos}}{Capacitancia de placas paralelas}{C = \frac{Q}{V}\\[2pt]C = \frac{\varepsilon_0\varepsilon_r A}{d}}
\formula{}{Capacitores en serie y en paralelo}{\text{Paralelo: } C_{eq} = C_1 + C_2\\[2pt]\text{Serie: } \frac{1}{C_{eq}} = \frac{1}{C_1} + \frac{1}{C_2}}
\formula{}{Capacitor con capas en serie}{C = \frac{\varepsilon_0 A}{t/\varepsilon_r + (d - t)}\\[2pt]E_{aire} = \varepsilon_r E_{diel}}
\formula{}{Capacitor con capas en paralelo}{C = \frac{\varepsilon_0}{d}\left(\varepsilon_{r1}A_1 + \varepsilon_{r2}A_2\right)\\[2pt]E_1 = E_2 = \frac{V}{d}}
\formula{}{Energía en el campo eléctrico}{W = \tfrac12 CV^2 = \int_v \tfrac12\varepsilon E^2\,dv\\[2pt]u_e = \tfrac12\varepsilon E^2}
\formula{}{Polarización, D y carga ligada}{\vec D = \varepsilon_0\vec E + \vec P\\[2pt]\vec P = \chi_E\varepsilon_0\vec E\\[2pt]\varepsilon_r = 1 + \chi_E\\[2pt]\rho_{s,b} = P = (\varepsilon_r - 1)\varepsilon_0E}
\formula{}{Condiciones de frontera entre dieléctricos}{E_{tan1} = E_{tan2}\\[2pt]D_{n1} - D_{n2} = \rho_S \quad (\rho_S = 0 \Rightarrow D_{n1} = D_{n2})}
\formula{}{Ley de Gauss para $\vec D$}{\oint_S \vec D\cdot d\vec s \;=\; Q_{\text{libre}}}
\formula{}{Capacitancia de un coaxial y de un capacitor esférico}{C_{coax} = \frac{2\pi\varepsilon L}{\ln(b/a)}\\[2pt]C_{esf} = \frac{4\pi\varepsilon}{1/a - 1/b}\\[2pt]RC = \frac{\varepsilon}{\sigma}}
\par\begin{minipage}{\linewidth}\grupo{Constantes}
\begin{tabular}{@{}l@{\ \ }l@{}}$\varepsilon_0$ & $8{,}854\times10^{-12}\ \mathrm{F/m} \approx \dfrac{10^{-9}}{36\pi}$\\[3pt]$\mu_0$ & $4\pi\times10^{-7}\ \mathrm{H/m}$\\[3pt]$q_e$ & $1{,}602\times10^{-19}\ \mathrm{C}$\\[3pt]$n_{Cu}$ & $8{,}5\times10^{28}\ \text{electrones}/\mathrm{m^3}$\\[3pt]$\sigma_{Cu},\ \sigma_{Al},\ \sigma_{Ag}$ & $5{,}8\,;\ 3{,}72\,;\ 6{,}17\ \times10^{7}\ \mathrm{S/m}$\\[3pt]$E_{rup,\,aire}$ & $\approx 3\ \mathrm{MV/m}$\end{tabular}\end{minipage}
\end{multicols}
\end{document}
